% Source: https://artofproblemsolving.com/community/c3h2319666_a_lot_of_not_very_hard_cp_problems_keeping_yall_busy
% Stable identity: aops-topic-2319666-q03
Two circles meet in at most $2$ points, two lines in at most $1$ point, and a line and a circle in at most $2$ points. Counting pairs, the number of intersection points is at most
\[2\dbinom52+1\cdot\dbinom42+2\cdot5\cdot4=20+6+40=\boxed{66}.\]
The figure shows the same count for $2$ circles and $2$ lines: $2+1+8=11$ points.

\begin{asy}
usepackage("amsmath");
size(5.5cm);
defaultpen(fontsize(10pt));
pair c1=(0,0), c2=(1.3,0.3); real r1=1.4, r2=1.1;
path C1=circle(c1,r1), C2=circle(c2,r2);
path L1=(-2.2,-0.9)--(3.2,0.99), L2=(-1.8,2.1)--(2.9,-1.425);
draw(C1,linewidth(0.9)); draw(C2,linewidth(0.9)); draw(L1,blue); draw(L2,red);
pair[] P;
P.append(intersectionpoints(C1,C2)); P.append(intersectionpoints(L1,L2));
P.append(intersectionpoints(C1,L1)); P.append(intersectionpoints(C1,L2));
P.append(intersectionpoints(C2,L1)); P.append(intersectionpoints(C2,L2));
for(int i=0;i<P.length;++i) dot(P[i],linewidth(4));
\end{asy}

The bound $66$ is attained. Take five circles of radius $1$ whose centres are distinct points on a short segment; any two of them meet in two points. Choose the centres so that no three circles pass through a common point (possible, since each of the finitely many intersection points rules out only finitely many positions for a further centre). Then draw four lines through the region common to all five discs, so each line crosses each circle twice; choose them with different slopes, no three through a point, and avoiding all previously found intersection points (again only finitely many directions or positions are excluded each time). Then all $66$ intersection points are distinct.
